\chapter{Fichiers de Configuration}
\label{annexe:configs}

Cette annexe regroupe les fichiers de configuration des différents composants de l'infrastructure.

\section{Configuration MikroTik CHR}
\label{sec:config_mikrotik}

Configuration complète du routeur MikroTik CHR : bridges par zone (VLAN), adresses IP des passerelles, règles de firewall inter-VLANs, et port mirroring vers Suricata.

\begin{lstlisting}[language=bash, caption={mikrotik-config.rsc --- Configuration du routeur MikroTik}, label={lst:mikrotik}]
# PFE IoT Security TT -- MikroTik CHR Configuration

/system identity set name=mikrotik-pfe

# Creer les bridges par VLAN
/interface bridge
add name=bridge-iot     comment="Zone IoT - VLAN 10"
add name=bridge-dmz     comment="Zone DMZ - VLAN 20"
add name=bridge-pki     comment="Zone PKI - VLAN 30"
add name=bridge-siem    comment="Zone SIEM - VLAN 40"
add name=bridge-analyse comment="Zone Analyse - VLAN 50"

# Assigner les ports aux bridges
/interface bridge port
add bridge=bridge-iot     interface=ether2
add bridge=bridge-dmz     interface=ether3
add bridge=bridge-pki     interface=ether4
add bridge=bridge-siem    interface=ether5
add bridge=bridge-analyse interface=ether6

# Adresses IP des passerelles (routeur L3)
/ip address
add address=192.168.1.2/24    interface=ether1       comment="Uplink pfSense"
add address=192.168.10.1/24   interface=bridge-iot    comment="GW Zone IoT"
add address=192.168.20.1/24   interface=bridge-dmz    comment="GW Zone DMZ"
add address=192.168.30.1/24   interface=bridge-pki    comment="GW Zone PKI"
add address=192.168.40.1/24   interface=bridge-siem   comment="GW Zone SIEM"
add address=192.168.50.1/24   interface=bridge-analyse comment="GW Zone Analyse"

# Route par defaut vers pfSense
/ip route
add dst-address=0.0.0.0/0 gateway=192.168.1.1

# Firewall -- ACL inter-VLANs
/ip firewall filter
add chain=forward connection-state=established,related action=accept
add chain=forward src-address=192.168.10.0/24 dst-address=192.168.20.10 \
    protocol=tcp dst-port=8883 action=accept comment="IoT->MQTT TLS"
add chain=forward src-address=192.168.10.0/24 dst-address=192.168.30.10 \
    protocol=tcp dst-port=8200 action=accept comment="IoT->Vault PKI"
add chain=forward src-address=192.168.20.0/24 dst-address=192.168.30.10 \
    protocol=tcp dst-port=8200 action=accept comment="DMZ->Vault verify"
add chain=forward src-address=192.168.20.0/24 dst-address=192.168.40.10 \
    protocol=tcp dst-port=1514 action=accept comment="DMZ->Wazuh logs"
add chain=forward src-address=192.168.50.0/24 dst-address=192.168.40.10 \
    protocol=tcp dst-port=1514 action=accept comment="Analyse->Wazuh"
add chain=forward src-address=192.168.10.0/24 dst-address=192.168.40.0/24 \
    action=drop comment="BLOCK IoT->SIEM direct"
add chain=forward src-address=192.168.10.0/24 dst-address=192.168.50.0/24 \
    action=drop comment="BLOCK IoT->Analyse direct"
add chain=forward action=drop comment="DROP all other inter-VLAN"
\end{lstlisting}

\section{Configuration Mosquitto (TLS/mTLS)}
\label{sec:config_mosquitto}

\begin{lstlisting}[language=bash, caption={mosquitto.conf --- Broker MQTT avec TLS 1.3 et mTLS}, label={lst:mosquitto_conf}]
# ==============================================
# Mosquitto v2.0 -- Configuration TLS 1.3 / mTLS
# PFE IoT Security TT
# ==============================================

# Listener TLS (port securise)
listener 8883

# Certificats serveur
cafile      /mosquitto/certs/ca-chain.crt
certfile    /mosquitto/certs/mosquitto.crt
keyfile     /mosquitto/certs/mosquitto.key

# mTLS : exiger certificat client
require_certificate true
use_identity_as_username true

# Version TLS minimale
tls_version tlsv1.3

# Ciphersuites TLS 1.3
ciphers_tls1.3 TLS_AES_256_GCM_SHA384:TLS_CHACHA20_POLY1305_SHA256

# ACL basee sur le CN du certificat
acl_file /mosquitto/config/aclfile

# Logging
log_dest stdout
log_type all
connection_messages true

# Securite
allow_anonymous false
max_connections 200
max_inflight_messages 20
max_queued_messages 1000

# Persistence
persistence true
persistence_location /mosquitto/data/
\end{lstlisting}

\section{Configuration Suricata (MQTT natif)}
\label{sec:config_suricata}

Extrait de la configuration Suricata v7.0 pertinente pour la détection MQTT.

\begin{lstlisting}[language=yaml, caption={suricata.yaml --- Extrait de la configuration IDS}, label={lst:suricata_yaml}]
# Suricata v7.0 -- Configuration pour detection MQTT
# PFE IoT Security TT

vars:
  address-groups:
    IOT_NET:     "192.168.10.0/24"
    DMZ_NET:     "192.168.20.0/24"
    MQTT_BROKER: "192.168.20.10"

app-layer:
  protocols:
    mqtt:
      enabled: yes
      ports: [1883, 8883]
    tls:
      enabled: yes
      ports: [8883, 443]

outputs:
  - eve-log:
      enabled: yes
      filetype: regular
      filename: eve.json
      types:
        - alert:
            tagged-packets: yes
            metadata: yes
        - mqtt:
            passwords: no
        - tls:
            extended: yes
        - stats:
            totals: yes
            threads: no

rule-files:
  - suricata.rules
  - mqtt-custom.rules
\end{lstlisting}

\section{Règles Suricata personnalisées}
\label{sec:config_suricata_rules}

\begin{lstlisting}[language=bash, caption={mqtt-custom.rules --- Règles IDS personnalisées MQTT}, label={lst:suricata_rules}]
# ===================================================
# Regles Suricata personnalisees -- Detection MQTT IoT
# PFE IoT Security TT
# ===================================================

# SID 1000001 -- Connexion MQTT sans TLS (port 1883)
alert mqtt any any -> $MQTT_BROKER 1883 (msg:"PFE - MQTT sans TLS detecte"; \
  flow:to_server,established; sid:1000001; rev:1; \
  classtype:policy-violation;)

# SID 1000002 -- Flood MQTT (>50 PUBLISH en 10s)
alert mqtt $IOT_NET any -> $MQTT_BROKER any (msg:"PFE - MQTT flood detecte"; \
  flow:to_server,established; mqtt.type:PUBLISH; \
  threshold:type both, track by_src, count 50, seconds 10; \
  sid:1000002; rev:1; classtype:attempted-dos;)

# SID 1000003 -- Subscribe wildcard '#'
alert mqtt $IOT_NET any -> $MQTT_BROKER any (msg:"PFE - MQTT subscribe wildcard #"; \
  flow:to_server,established; mqtt.type:SUBSCRIBE; \
  content:"#"; sid:1000003; rev:1; classtype:policy-violation;)

# SID 1000004 -- Payload anormalement grand (>4096 octets)
alert mqtt $IOT_NET any -> $MQTT_BROKER any (msg:"PFE - MQTT payload > 4096B"; \
  flow:to_server,established; mqtt.type:PUBLISH; \
  dsize:>4096; sid:1000004; rev:1; classtype:bad-unknown;)

# SID 1000005 -- Tentative connexion depuis zone non-IoT
alert mqtt !$IOT_NET any -> $MQTT_BROKER 8883 (msg:"PFE - MQTT depuis zone non-IoT"; \
  flow:to_server; sid:1000005; rev:1; classtype:policy-violation;)
\end{lstlisting}

\section{Règles Wazuh personnalisées}
\label{sec:config_wazuh_rules}

\begin{lstlisting}[language=xml, caption={local\_rules.xml --- Règles Wazuh pour corrélation MQTT/Suricata}, label={lst:wazuh_rules}]
<!-- Wazuh custom rules -- PFE IoT Security TT -->
<group name="pfe-iot,suricata,mqtt">

  <!-- Alerte Suricata MQTT generique -->
  <rule id="100100" level="5">
    <if_sid>86601</if_sid>
    <field name="alert.signature_id">^1000</field>
    <description>PFE -- Alerte Suricata MQTT detectee</description>
    <group>pfe-iot,mqtt,suricata</group>
  </rule>

  <!-- MQTT flood (DoS) -->
  <rule id="100101" level="10">
    <if_sid>100100</if_sid>
    <field name="alert.signature_id">1000002</field>
    <description>PFE -- MQTT DoS flood detecte (>50 msg/10s)</description>
    <group>pfe-iot,mqtt,dos</group>
    <options>alert_by_email</options>
  </rule>

  <!-- Subscribe wildcard -->
  <rule id="100102" level="8">
    <if_sid>100100</if_sid>
    <field name="alert.signature_id">1000003</field>
    <description>PFE -- MQTT subscribe wildcard # detecte</description>
    <group>pfe-iot,mqtt,policy-violation</group>
  </rule>

  <!-- Connexion MQTT sans TLS -->
  <rule id="100103" level="12">
    <if_sid>100100</if_sid>
    <field name="alert.signature_id">1000001</field>
    <description>PFE -- MQTT sans TLS (violation de politique)</description>
    <group>pfe-iot,mqtt,tls-violation</group>
    <options>alert_by_email</options>
  </rule>

  <!-- Correlation : 3+ alertes MQTT du meme device en 60s -->
  <rule id="100110" level="13" frequency="3" timeframe="60">
    <if_matched_sid>100100</if_matched_sid>
    <same_source_ip />
    <description>PFE -- Activite suspecte repetee depuis le meme device IoT</description>
    <group>pfe-iot,mqtt,correlation</group>
  </rule>

</group>
\end{lstlisting}

\section{Dockerfile du simulateur de flotte}
\label{sec:config_dockerfile}

\begin{lstlisting}[language=bash, caption={Dockerfile --- Image Docker du simulateur de flotte IoT}, label={lst:dockerfile_fleet}]
FROM python:3.11-alpine

RUN apk add --no-cache ca-certificates openssl \
    && update-ca-certificates

WORKDIR /app
COPY app/requirements.txt /app/requirements.txt
RUN pip install --no-cache-dir -r /app/requirements.txt

COPY app/fleet_sim.py /app/fleet_sim.py
COPY certs/ /app/certs/
COPY dataset/ /app/dataset/

CMD ["python", "/app/fleet_sim.py", \
     "--csv", "/app/dataset/iot_communication_security_dataset_sample.csv"]
\end{lstlisting}

\section{Dockerfile Toolbox}
\label{sec:config_dockerfile_toolbox}

\begin{lstlisting}[language=bash, caption={Dockerfile --- Image toolbox pour tests réseau et sécurité}, label={lst:dockerfile_toolbox}]
FROM alpine:3.20

RUN apk add --no-cache \
    bash ca-certificates curl openssl \
    bind-tools iproute2 iputils tcpdump \
    net-tools nmap jq \
    python3 py3-pip git \
    busybox-extras tzdata \
  && update-ca-certificates

RUN adduser -D pfe && mkdir -p /work && chown -R pfe:pfe /work
WORKDIR /work
USER pfe

CMD ["bash"]
\end{lstlisting}

\section{Commandes de création des réseaux Docker}
\label{sec:config_docker_networks}

\begin{lstlisting}[language=bash, caption={Création des réseaux Docker par zone de sécurité}, label={lst:docker_networks}]
# VLAN 10 -- Zone IoT
docker network create --driver bridge \
  --subnet 192.168.10.0/24 --gateway 192.168.10.1 \
  --opt com.docker.network.bridge.name=br-iot pfe-iot-network

# VLAN 20 -- Zone DMZ
docker network create --driver bridge \
  --subnet 192.168.20.0/24 --gateway 192.168.20.1 \
  --opt com.docker.network.bridge.name=br-dmz pfe-dmz-network

# VLAN 30 -- Zone PKI
docker network create --driver bridge \
  --subnet 192.168.30.0/24 --gateway 192.168.30.1 \
  --opt com.docker.network.bridge.name=br-pki pfe-pki-network

# VLAN 40 -- Zone SIEM
docker network create --driver bridge \
  --subnet 192.168.40.0/24 --gateway 192.168.40.1 \
  --opt com.docker.network.bridge.name=br-siem pfe-siem-network

# VLAN 50 -- Zone Analyse
docker network create --driver bridge \
  --subnet 192.168.50.0/24 --gateway 192.168.50.1 \
  --opt com.docker.network.bridge.name=br-analyse pfe-analyse-network
\end{lstlisting}

\section{Variables d'environnement du projet}
\label{sec:config_env}

\begin{table}[H]
\centering
\caption{Variables d'environnement principales du projet}
\label{tab:env_vars}
\begin{tabular}{|l|l|p{6cm}|}
\hline
\textbf{Variable} & \textbf{Valeur} & \textbf{Description} \\
\hline
\texttt{VAULT\_ADDR} & \texttt{http://192.168.30.10:8200} & Adresse de l'API Vault \\ \hline
\texttt{VAULT\_TOKEN} & \texttt{<root\_token>} & Token d'authentification \\ \hline
\texttt{MQTT\_BROKER} & \texttt{192.168.20.10} & IP du broker Mosquitto \\ \hline
\texttt{MQTT\_PORT\_TLS} & \texttt{8883} & Port MQTT sécurisé \\ \hline
\texttt{MQTT\_PORT\_TCP} & \texttt{1883} & Port MQTT non sécurisé \\ \hline
\texttt{CA\_CERT} & \texttt{/certs/ca-chain.crt} & Chaîne CA (root + intermediate) \\ \hline
\texttt{WAZUH\_MANAGER} & \texttt{192.168.40.10} & IP du manager Wazuh \\ \hline
\texttt{SURICATA\_EVE} & \texttt{/var/log/suricata/eve.json} & Fichier événements Suricata \\ \hline
\end{tabular}
\end{table}
